% PLAN:N7
% NODES: res:sensing-grade-structural res:window-verdict irf:n14-window-chi-flat-through-n20 thm:readability-hardness-nogo irf:local-trades-absent-generic irf:denominators-grow-with-n irf:n14b-delocalization-only-far-above-ceiling def:localization-protection irf:stoquastic-cure-check irf:splitting-vs-clock def:planted-block thm:planted-escape res:planted-trade-family irf:xeb-epistemology-at-exponential-g fork:planted-hardness-formalisation
\section{Limits and discussion}\label{sec:N7}

% N7-p1
Within the validated operating window the tested generic distinct-weight instances are entangling but localized - the required bond dimension sits in $\{2, 3, 4\}$ and stays flat from $n = 8$ to $20$, and delocalization appears only far above the fingerprint ceiling. The benchmark therefore delivers a working spectral-sensing certificate on the generic family, its value standing on the fingerprint and decoder legs. The discussion separates three results the paper establishes at different strengths. The general acquisition-reconstruction architecture is exact and target-independent. The benchmark-specific diagnostics are measured over the stated window. The exploratory complexity extension of the planted family opens a question rather than settling one. The architecture is the widely applicable result. It asks only for a single coherent propagator expectation $G(t) = \langle\psi|e^{-iHt}|\psi\rangle$ - the quantity a universal quantum simulator produces \cite{lloyd1996universal, georgescu2014sim} and the real-time spectral estimators already sample \cite{somma2019eigenvalue, obrien2019qpe}. Any platform or protocol that can furnish that overlap therefore inherits the whole surface-and-response construction, with the mathematics carrying every later step on the banked record \cite{mccaul2025snacks, leamer2024emulation}. This section is where every global complexity boundary is stated in full.

% N7-p2
That flat band asks for an explanation, and we offer one as a finite-size compressibility proposition, hypotheses explicit. Take $H = H_{0} + JV$ on the generic distinct-weight family, with $V$ a bounded noncommuting perturbation of unit per-bond strength and $J$ inside the fingerprint window. Assume first the absence of local signed coincidences relevant to $V$, under which the off-resonance gap obeys $D_{\min} = \Omega(n)$ with probability $1 - \delta$ over the production ensemble. Assume second the local perturbation norm. Assume third a prethermal window $t^{*} \sim e^{c' n / J}$ covering the record horizon $t_{\max}(\tau)$~\eqref{res:aperture-resolution} \cite{abanin2017prethermal}. A single off-resonant rotation then makes the leading effective Hamiltonian diagonal, with
\begin{widetext}
\begin{equation*}
\begin{aligned}
&\lVert S^{(i)}\rVert \le \frac{J}{D_{\min}}\,, \qquad
\lVert R^{(i)}\rVert \le \frac{c\,J^{2}}{D_{\min}} = O\!\left(\frac{J^{2}}{n}\right), \qquad
\lVert R\rVert \sim \frac{n J^{2}}{D_{\min}} \sim J^{2}\,,\\
&\chi_{\varepsilon} = 2^{O(1)}\,\operatorname{poly}(1/\varepsilon) \quad \text{at fixed accuracy } \varepsilon\,, \ \text{independent of } n\,,
\end{aligned}
\end{equation*}
\end{widetext}
so the dressed residual is weak per bond while its global norm is a constant that does not vanish, and within the controlled window the physical state admits an efficient approximate matrix product state \cite{verstraete2008mps}. The proposition sits with the locality-bound and short-time literature, in which finite-window dynamics of local Hamiltonians is classically representable \cite{lieb1972finite, wild2023shorttime}. It is a finite-window compressibility statement about this construction, and it speaks to the tested range rather than to all classical representations or the worst case. The $e^{-S}$ undressing generally adds entanglement, so the physical state carries no exact bond-dimension-two guarantee.

% N7-p3
The quantitative engine of the proposition is combinatorial. On generic distinct-weight instances local trades - pairs of computational configurations differing on a bounded site set with equal instance energy - are absent, so every off-diagonal matrix element of $V$ is off-resonant by at least the local weight mismatch. The transition denominators are signed weight-sums over the move's window, $D_{xy} = \sum_{i} \epsilon_{i} a_{i}$. Distinctness alone secures $D_{xy} \neq 0$ for the two-site moves of $V_{\mathrm{XXZ}}$, where $D_{xy} \in \{\pm(a_{i} + a_{i+1}), \pm(a_{i} - a_{i+1})\}$, while a general $k$-local $V$ requires the explicit hypothesis that local signed coincidences relevant to $V$ are absent. Under that hypothesis, with the $a_{i}$ drawn i.i.d.\ from $\mathrm{Uniform}\{1, \dots, n^{2}\}$, a first-moment bound over the $O(n)$ nearest-neighbour differences gives $\Pr(D_{\min} \le \gamma) \le 2(n-1)\gamma/n^{2}$, so $D_{\min} \ge (\delta/2)\,n$ with probability $1 - \delta$ and the gap grows linearly in $n$. The bound is probabilistic over the production ensemble rather than a per-instance certainty.

% N7-p4
The other side of the coupling axis is measured rather than argued. The fingerprint ceiling is absolute - the filter window cannot pass couplings beyond roughly $J = 0.1$ to $0.2$ without merging unit-spaced clusters, and no $\tau$ setting rescues it - while delocalization requires $J$ at the local-resonance scale. Measured, $J = 1$ at ten times the ceiling gives bond dimensions $4/4/5$ across the sizes measured, a marginal onset, and $J = 3$ at thirty times the ceiling gives $6/6/11$ with the half-cut entropy $S_{\mathrm{vN}}$ rising from $0.96$ to $1.42$, so there the bond dimension grows with $n$. Entangling growth exists in this model, at couplings far above where the exact read can live, and the protocol's readable window does not reach those couplings.

% N7-p5
One mechanism accounts for both sides. The operating model is a random-field XXZ chain whose diagonal disorder comes from the instance weights,
\begin{equation*}
W \sim n^{2}\,, \qquad \frac{W}{J} \sim \frac{n^{2}}{J} \quad \text{at fixed in-window } J\,,
\end{equation*}
so at fixed coupling the disorder outruns the perturbation as $n$ grows. Localization deepens with size and in-window entanglement stays bounded - the flat bond dimension from $n = 8$ to $20$ - while raising $J$ towards the local-resonance scale closes the separation and produces the growth measured far above the ceiling. The picture is the established one for random-field XXZ chains: the numerically mapped localization of the random-field magnet \cite{znidaric2008mbl}, its phase transition \cite{pal2010mbl}, its description through local conservation laws \cite{serbyn2013liom} and its rigorous construction for spin chains \cite{imbrie2016mbl}. Because the disorder scale grows with $n$ by construction the family also sits near the strongly tilted systems studied as Stark localization, where fields that grow with system size localize without fixed-strength randomness \cite{schulz2019stark}, the localization a growing field induces in clean interacting chains \cite{vannieuwenburg2019stark} as observed without disorder \cite{morong2021stark}. Growing disorder is a regime the fixed-strength localization theorems do not directly cover, so we offer the mechanism as the explanation of the tested range, not as an asymptotic theorem about the family.

% N7-p6
The localization literature itself bounds what the mechanism can promise. The asymptotic status of many-body localization in one dimension is contested \cite{nandkishore2015mbl, abanin2019mbl}. Avalanches seeded by rare thermal regions can destabilise localization at large sizes and long times \cite{deroeck2017avalanche, luitz2017thermalize}, finite-size numerics are known to drift towards the thermal side as sizes grow \cite{morningstar2022avalanches}, and even deep in localized regimes entanglement grows logarithmically in time without bound \cite{bardarson2012unbounded}. Each of these failure modes concerns scales beyond the tested range, and none is observed within it. The caveats bound extrapolation beyond $n = 20$ and beyond the record horizon - they do not weaken the reported measurements.

% N7-p7
The family invites a check against the sign-problem literature, made here explicitly. The chain is one-dimensional, nearest-neighbour and bipartite - conjugating by $Z$ on the odd sublattice sends $XX + YY$ to $-(XX + YY)$ and leaves $ZZ$ and the diagonal $H_{0}$ invariant - so $H = H_{0} + J V_{\mathrm{XXZ}}$ is stoquastic up to a local basis change in the standard taxonomy \cite{bravyi2008stoquastic, klassen2019stoquastic}. Quantum Monte Carlo therefore faces no thermal sign problem in principle on this family \cite{troyer2005computational}, and because deciding whether a Hamiltonian can be cured of its sign problem is in general hard \cite{marvian2019cure, klassen2019hardness} we exhibit the transformation rather than assume it, in the spirit of methods that ease the sign problem where it cannot be removed \cite{hangleiter2020easing}. The surviving classical wall is observable access - the protocol quantities are filtered real-time-record objects, and extracting them from imaginary-time data is the ill-posed analytic-continuation problem, weaker than a sign problem and a statement about access to real-time observables rather than about the existence of an efficient equilibrium sampler.

% N7-p8
The generic bounds separate into two different scales. Readability constrains the instantaneous splitting scale $J\,B(L)$ - the staircase must stay inside the filter window at every size. Dynamics, and the entanglement it generates, depends on the integrated clock $Jt$. These are different quantities. A family that holds $J\,B(L)$ fixed while letting $Jt$ grow can therefore retain readability while its classical simulation cost grows, without contradicting the compressibility of the generic family. The pivot by itself produces no growth - the construction and the measurements that follow carry that burden.

% N7-p9
The escape family plants an equal-weight block of $L$ sites in the instance,
\begin{equation*}
a_{i} \equiv \bar{a} \ \text{ on the block}\,, \qquad H_{0}\big|_{\mathcal{M}} = E_{\mathcal{M}}\,, \qquad J = \frac{c}{L}\,,
\end{equation*}
where $\mathcal{M}$ is the block's degenerate manifold with common diagonal energy $E_{\mathcal{M}}$. Inside $\mathcal{M}$ the diagonal disorder cancels exactly, so $J$ acts there as a pure clock. The choice $J = c/L$ holds the worst in-window staircase spread bounded at every $n$ - the readability half - while entanglement accumulates on the integrated clock $Jt$, with measured bond dimension close to $2^{n/4}$. The construction follows the planted-instance tradition of hiding designed structure in an otherwise generic instance \cite{barthel2002planted}, and the degenerate manifold in which the clock runs is a designed analogue of the dynamically decoupled subspaces studied as scars and as Hilbert-space fragmentation \cite{turner2018scars, sala2020fragmentation}. The cancellation is exact only inside the planted manifold - outside it the generic disorder structure persists - and $J = c/L$ is a design choice of the family rather than a derived necessity.

% N7-p10
The theorem that results states the coexistence with both halves explicit,
\begin{widetext}
\begin{equation*}
J = \frac{c}{L}\,, \qquad J\,B(L) = 2JL = 2c \ \ \text{at every } n\,, \qquad Jt \ \ \text{growing on the protocol's own } \tau \text{ dial}\,.
\end{equation*}
\end{widetext}
With the block planted, in-manifold disorder cancels exactly and $J$ is a pure clock. The choice $J = c/L$ keeps the worst in-window spread bounded, so readability holds at every $n$, while the evolution time available to the protocol stretches polynomially on its own $\tau$ dial, so the integrated clock $Jt$ grows without breaching the window. The coexistence absent in the tested generic family - readability together with growing classical simulation cost - is realised by design. The theorem establishes that coexistence for the designed family over the analysed window, and it does not upgrade the generic verdict.

% N7-p11
Both halves of the trade carry numbers. Exact diagonalisation at $n = 12$, $L = 6$, $J = 0.033$ gives worst dressed in-window spread $0.413$, against both the $0.5$ threshold and the leading estimate $2JL = 0.4$, with the dressing shrinking as $O(1/n)$. TEBD at identical $(n, J, t)$ gives planted bond dimension $4/8/16/30$ at $n = 8/12/16/20$, close to $2^{n/4}$ with no cap hits, while the generic control stays pinned at bond dimension $2$ \cite{verstraete2008mps}. The clock cost is accounted through the acquisition ledger by the single substitution $t_{\max} = O(L)$ in the aperture reach~\eqref{res:aperture-resolution}: the supported filter-width domain widens to $\tau_{\max} = O(L^{2}/\log(1/\epsilon_{w}))$, the retained-node count grows as $\bar{m} = O(\lVert H\rVert L)$, the total shot budget as $S_{\mathrm{tot}} = O(\lVert H\rVert L)\,K$, and the per-window visibility floor stays flat in $L$ under a centred probe - the instantaneous spectral width is exchanged for a linearly longer clock at these derived acquisition costs, not obtained for free. Figure~\ref{fig:nf7} shows the trade: the worst in-window spread against $n$ with the $0.5$ threshold in panel (a), the required bond dimension for the planted family and the generic control with the $2^{n/4}$ guide in panel (b), the half-cut entanglement entropy against the accumulated clock $Jt$ in panel (c). The growth is consistent with exponential over the sizes measured. Panel (a) carries exact-diagonalisation points at $n = 8$ and $12$ only, exact diagonalisation being infeasible at $16$ and $20$, so readability at the larger sizes rests on the design bound rather than direct verification.

% N7-p12
At exponential degeneracy $g$ the per-step exact-integer read gives way to statistical certification against an exact prediction - the register of cross-entropy benchmarking \cite{boixo2018xeb}. The reference stays classically computable, because the planted block's in-manifold spectrum is exact by construction. The certificate changes epistemic register, not existence. That change bounds what any sampling-style comparison on the planted family could certify, and is part of the family's cost.

% N7-p13
The architecture the paper establishes is a reusable spectral-acquisition record and the family of reconstructions it supports, validated on the controlled benchmark across its measured operating window. Separately, the planted family opens a complexity question: whether the overlap record $A(t)$ on planted instances is formally hard to compute classically - the one question between the demonstrated family and a quantum-advantage-relevant claim. A single planted block is integrable-adjacent, its in-manifold dynamics close to a free precession, so the natural route is multi-block or junction constructions that harden the manifold, and we leave that formalisation as a named open problem. The best classical subset-sum algorithms set the benchmark any such claim would face \cite{bonnetain2020improved}. The observed bond-dimension growth is the cost of one simulation method over the analysed range, and whether it reflects a genuine classical lower bound is exactly what that open problem asks.
